% ============================================================
% 4_F_compatibility_test.tex — 兼容性测试
% 编写人：F（D代理完成）  日期：2026-07-22
% ============================================================

\section{兼容性测试}

\subsection{浏览器兼容性}

\begin{table}[htbp]
    \centering
    \caption{浏览器兼容性测试结果}
    \begin{tabular}{|c|c|c|c|c|}
        \hline
        \textbf{浏览器} & \textbf{版本} & \textbf{核心功能} & \textbf{样式渲染} & \textbf{结论} \\
        \hline
        Google Chrome & 120+ & 正常 & 正常 & 通过 \\
        \hline
        Microsoft Edge & 120+ & 正常 & 正常 & 通过 \\
        \hline
        Mozilla Firefox & 120+ & 正常 & 富文本编辑器光标偶有偏移 & 通过（轻微） \\
        \hline
        Apple Safari & 17+ & 正常 & 部分CSS Grid间距偏差1-2px & 通过（轻微） \\
        \hline
    \end{tabular}
\end{table}

\subsection{后端接口兼容性}

\begin{table}[htbp]
    \centering
    \caption{API响应格式一致性验证}
    \begin{tabular}{|c|c|c|}
        \hline
        \textbf{验证项} & \textbf{标准} & \textbf{结果} \\
        \hline
        成功响应格式 & \texttt{\{"code":200,"message":"success","data":...\}} & 全部一致 \\
        \hline
        错误响应格式 & \texttt{\{"code":xxx,"message":"...","data":null\}} & 全部一致 \\
        \hline
        分页响应格式 & \texttt{\{"content":[...],"totalElements":N,...\}} & 全部一致 \\
        \hline
        Content-Type & application/json;charset=UTF-8 & 全部一致 \\
        \hline
        空List返回 & \texttt{[]} 而非 null & 全部一致 \\
        \hline
    \end{tabular}
\end{table}

\subsection{数据库兼容性}

\begin{table}[htbp]
    \centering
    \caption{数据库兼容性验证}
    \begin{tabular}{|c|c|c|}
        \hline
        \textbf{数据库} & \textbf{测试方式} & \textbf{结果} \\
        \hline
        PostgreSQL 17 (生产) & Docker Compose 部署 & 全部功能正常 \\
        \hline
        H2 (开发) & dev profile 启动 & 全部功能正常 \\
        \hline
        数据迁移一致性 & 同表结构H2→PostgreSQL & 兼容，通过 \\
        \hline
    \end{tabular}
\end{table}

\subsection{操作系统兼容性}

\begin{table}[htbp]
    \centering
    \caption{操作系统兼容性}
    \begin{tabular}{|c|c|c|}
        \hline
        \textbf{操作系统} & \textbf{部署方式} & \textbf{结果} \\
        \hline
        Windows 11 & Docker Desktop (WSL2) & 全部正常 \\
        \hline
        Windows 11 & 本地 mvnw + npm run dev & 全部正常 \\
        \hline
    \end{tabular}
\end{table}

\subsection{兼容性测试结论}

系统在四大主流浏览器、两种部署模式（Docker/本地）、两种数据库（PostgreSQL/H2）下均表现一致，兼容性良好。Firefox 和 Safari 的轻微样式差异不影响功能使用，属于可接受范围。
